% !PRINTED 38
\weberpage{38}
%<<CONT

also, wenn mit $p,p'$ ganze rationale Functionen bezeichnet werden:
\begin{equation}
 A'''=pA+p'A'.
\tag{21}
\end{equation}
Setzt man die Ausdrücke (20), (21) in die dritte Gleichung (3) ein, so ergiebt sich für $A''''$ wieder ein Ausdruck von der Form (21) und so kann man fortfahren, und erhält schliesslich:
\begin{equation}
 A^{(\nu)}=PA+P'A',
\tag{22}
\end{equation}
worin $P,P'$ ganze rationale Functionen sind, deren Coëfficienten durch rationale Rechenoperationen aus den Coëfficienten von $A$ und $A'$ zusammengesetzt sind.

In der Formel (22) ist $A^{(\nu)}$ eine Constante. Besonders wichtig ist dieser Satz in dem Falle, wo $A^{(\nu)}$ von Null verschieden, also $A,A'$ relativ prim sind. Setzen wir in diesem Falle
\[
 P=A^{(\nu)}F(x),\qquad P'=A^{(\nu)}\Phi(x),
\]
so können wir nach Weglassung des Factors $A^{(\nu)}$ dem Satz folgenden Ausdruck geben:

\gesperrt{I. Sind $f(x)$ und $\psi(x)$ zwei ganze Functionen ohne gemeinsamen Theiler, so kann man zwei andere ganze Functionen $F(x)$ und $\Phi(x)$ bestimmen, die der Gleichung}
\begin{equation}
 F(x)f(x)+\Phi(x)\psi(x)=1
\tag{23}
\end{equation}
\gesperrt{identisch genügen.}

Der Satz lässt sich noch verallgemeinern. Multipliciren wir die Gleichung (23) mit einer beliebigen ganzen Function $\chi(x)$, so folgt:
\begin{equation}
 F(x)\chi(x)f(x)+\Phi(x)\chi(x)\psi(x)=\chi(x),
\tag{24}
\end{equation}
und nach \S~3 können wir
\begin{equation}
 \Phi(x)\chi(x)=Q(x)f(x)+\varphi(x)
\tag{25}
\end{equation}
setzen, so dass $Q(x),\varphi(x)$ ganze rationale Functionen von $x$ sind und der Grad von $\varphi(x)$ kleiner ist als der von $f(x)$. Setzen wir dies in (24) ein und setzen an Stelle von $F(x)\chi(x)+Q(x)\psi(x)$ wieder $F(x)$, so erhalten wir:
\[
 F(x)f(x)+\varphi(x)\psi(x)=\chi(x).
\]
Wir können daher den vorigen Satz so verallgemeinern:

\gesperrt{II. Sind $f(x),\psi(x),\chi(x)$ gegebene ganze rationale Functionen, und $f(x)$ und $\psi(x)$ ohne gemeinsamen Theiler, so lassen sich die ganzen}
%CONT>>
